% !TEX program = pdflatex
%
% Estimation de probabilités sur tableaux d'élimination via marches quantiques
% et validation en direct sur marchés de prédiction
%
% Préprint arXiv, version française — Juin 2026
% Catégories : cs.ET (Emerging Technologies) ; cross-list : quant-ph, q-fin.CP

\documentclass[11pt]{article}
\usepackage[utf8]{inputenc}
\usepackage[T1]{fontenc}
% babel-french n'est pas systématiquement installé ; on protège l'inclusion
% par \IfFileExists pour rester portable. Les guillemets \og / \fg sont
% définis plus bas par \providecommand au cas où babel ne les fournit pas.
\IfFileExists{french.ldf}{\usepackage[french]{babel}}{}
\usepackage{lmodern}
\usepackage{microtype}
\usepackage[a4paper,
            top=1in, bottom=1in,
            left=1.15in, right=1.15in]{geometry}
\usepackage{amsmath, amssymb, amsthm}
\usepackage{graphicx}
\usepackage{hyperref}
\hypersetup{colorlinks=true, urlcolor=blue, citecolor=blue,
            linkcolor=black}
\usepackage{booktabs}
\usepackage{authblk}
\usepackage{enumitem}
\usepackage{xcolor}
\usepackage{placeins}

\usepackage{titlesec}
\titlespacing*{\section}{0pt}{2.0ex plus 0.3ex minus 0.2ex}{1.2ex plus 0.2ex}
\titlespacing*{\subsection}{0pt}{1.6ex plus 0.3ex minus 0.2ex}{0.8ex plus 0.2ex}

\widowpenalty=10000
\clubpenalty=10000

\setlength{\textfloatsep}{14pt plus 4pt minus 2pt}
\setlength{\intextsep}{12pt plus 3pt minus 2pt}

% Définitions sûres pour les guillemets français (en cas de babel-french trop
% ancien sans \og / \fg).
\providecommand{\og}{\guillemotleft\,}
\providecommand{\fg}{\,\guillemotright}

\title{Estimation de probabilités sur tableaux d'élimination\\
via marches quantiques et validation en direct\\
sur marchés de prédiction}

\author{Jean-Thomas Dauchel\thanks{\texttt{contact@pulsquant.com}}}

\affil{PulsQuant --- Recherche indépendante}

\date{Juin 2026}

\begin{document}
\maketitle

\begin{abstract}
Ce papier applique les marches quantiques aux tableaux d'élimination
des tournois sportifs. À notre connaissance c'est le premier
traitement formel du sujet : la note \emph{Quantum Frontiers} de 2018
\cite{quantumfrontiers_wc_2018} esquissait bien un encodage en qudits
pour la Coupe du Monde FIFA mais n'a jamais été soumise à publication,
et le récent papier QNN football de Nature/Scientific Reports
\cite{nature_qnn_football_2025} attaque une tâche différente
(classification match par match en QML).

On mappe un tableau à $n$ équipes à élimination directe sur un circuit
de $n{-}1$ qubits avec portes $R_Y$ contrôlées. Les matchs du premier
tour utilisent la logistique Elo exacte ; les tours suivants utilisent
une approximation \emph{graînée} (\og seeded \fg{}) dont on quantifie
le biais par rapport à une énumération en forme close suivant la même
règle d'appariement, de sorte que toute divergence $Q$--$C$ s'impute à
l'approximation et non à un artefact d'implémentation quantique.

L'encodage est testé rétroactivement sur cinq vrais tableaux
2025--2026 chargés depuis une base de marchés de prédiction en
production (pas des hypothèses équilibrées). Sur les Playoffs NBA
2026, où les San Antonio Spurs ont atteint les Finales NBA depuis la
8\textsuperscript{e} tête de série de l'Ouest en éliminant le numéro 1
OKC en finale de conférence Ouest, la marche graînée attribue
$P_Q(\text{SAS atteint les Finales}) = 5{,}87\,\%$ contre
$P_C = 3{,}43\,\%$ (amplification $\mathbf{1{,}71\times}$), et une
amplification $\mathbf{3{,}05\times}$ sur $P(\text{SAS champion})$.
Sur le tableau de Roland Garros femmes à partir du R16, la championne
réelle Mirra Andreeva voit son marginal de titre amplifié
$\mathbf{1{,}68\times}$, et la finaliste outsider Maja Chwalinska
(Tennis Elo 1710, classée 114\textsuperscript{e} WTA avant le tournoi)
voit sa probabilité de finale amplifiée $\mathbf{2{,}78\times}$. Les
demi-finalistes outsiders ATP (Jakub Mensik, Matteo Arnaldi)
ressortent à $\mathbf{4{,}17\times}$ et $\mathbf{4{,}44\times}$. Deux
contrôles négatifs propres confirment que la méthode n'est pas qu'un
amplificateur d'upsets : la victoire \og assurée \fg{} du PSG en
Ligue des Champions donne $Q/C \approx 1{,}01$, et Vegas sortant
Colorado en finale de conférence Ouest NHL (statistiquement un
pile-ou-face malgré le récit d'upset) donne $Q/C = 0{,}63\times$ sur
Vegas atteignant la finale.

Comme prolongements opérationnels de l'encodage, on implémente sur la
même requête (SAS champion sur le tableau NBA) à la fois
l'amplification d'amplitude de Grover et l'estimation d'amplitude
quantique \cite{brassard2002quantum}. La courbe d'amplification atteint
son pic à $97{,}89\,\%$ sur $k = 4$ itérations (optimum théorique du
manuel), et la QAE complète avec $m = 6$ qubits ancillaires (21 qubits
au total sur un simulateur Aer monoposte) renvoie
$\hat p_{\text{mode}} = 3{,}81\,\%$ sur l'amplitude SAS-champion
contre la référence $p_0 = 3{,}46\,\%$ obtenue par échantillonnage non
amplifié. Le mode QAE et l'estimateur en moyenne pondérée s'accordent
avec la référence à l'intérieur de la bande de résolution $\pi / 2^m$
pour tout $m \in \{4, 5, 6\}$, réalisant le speedup quadratique de
Brassard sur une requête de tableau réel.

Une prédiction prospective pour la Coupe du Monde FIFA 2026
(coup d'envoi 11 juin 2026) est émise dans la même forme : nations de
milieu de hiérarchie (NED, BEL, AUT, COL) amplifiées de 2 à 3 points
de pourcentage, favoris européens de tête (FRA, ESP, BRA) réduits
d'autant. ARG fait exception parmi les favoris de tête --- amplifiée
$+6{,}14$\,pp parce que son chemin dans le tableau équilibré traverse
un quart Conmebol plus faible.

À l'échelle simulateur (15 qubits), nous ne revendiquons aucun
avantage quantique. La contribution est conceptuelle : l'encodage
capture les corrélations structurelles du tableau que le Monte Carlo
classique ne reproduit qu'au prix d'un conditionnement explicite
chemin par chemin, et la construction passe à l'échelle de manière
tractable vers les 50 à 100 qubits logiques attendus du matériel
tolérant aux fautes de première génération.
\end{abstract}

\section{Introduction}
\label{sec:intro}

Les tournois sportifs à élimination directe présentent de fortes
\emph{corrélations structurelles} entre les résultats des matchs : si
une tête de série élevée perd tôt, les matchs en aval sont rebattus de
manière à modifier la probabilité marginale de titre de tous les
autres participants. Pricer correctement ces corrélations est un
problème central pour les marchés de prédiction et les bookmakers, qui
doivent agréger l'espérance sur un nombre exponentiel de chemins du
tableau pour émettre des cotes futures sur le titre. Le Monte Carlo
classique (MC) est l'approche standard
\cite{glasserman2004monte}, avec l'estimateur de manuel faisant tourner
$N$ simulations indépendantes du tableau complet et comptant les
issues championnes \cite{simulate_bracket_pulsquant}. Chaque simulation
est en $O(n)$ dans le nombre de matchs $n$, et le bruit
d'échantillonnage MC sur le marginal d'un champion à faible
probabilité passe à l'échelle en $O(1/\sqrt{N})$. Côté modélisation
match par match, les références classiques les plus proches sont la
famille Dixon--Coles des modèles football basés sur Poisson
\cite{dixon1997modelling} et leurs extensions bayésiennes
\cite{constantinou2017bayesian} : elles fixent les probabilités de
match implicites, tandis que la corrélation structurelle entre matchs
est typiquement laissée à la couche MC.

La saison 2025--2026 a livré un ensemble inhabituellement riche
d'upsets pour mettre un tel estimateur à l'épreuve. Aux Playoffs NBA,
les San Antonio Spurs ont atteint les Finales NBA depuis la
8\textsuperscript{e} tête de série de l'Ouest, en éliminant OKC Thunder
(tête de série numéro 1) en finale de conférence Ouest. Roland Garros
2026 a perdu trois des quatre têtes de série avant les quarts (Sinner
côté hommes, Świątek et Sabalenka côté femmes). En NHL, les Vegas
Golden Knights ont éliminé Colorado, meilleure équipe de la saison
régulière. Aucun de ces résultats n'est statistiquement extrême sur
son marginal de match isolé : un Elo bien calibré leur donne à chacun
une probabilité à deux chiffres. Ce qui les rend difficiles à pricer,
c'est le poids joint qu'ils portent dans l'arbre du tableau une fois
qu'on conditionne sur un parcours profond. La corrélation entre tours
successifs amplifie le marginal de l'outsider survivant bien au-delà
de ce que suggère une approximation tours indépendants.

Ce papier propose une marche quantique qui encode ces corrélations par
l'intrication entre qubits de matchs conditionnellement dépendants, et
en étudie le comportement dans trois cadres : un tableau synthétique à
8 équipes en sanity check ; cinq vrais tournois 2025--2026 backtestés
contre la vérité terrain (Playoffs NBA, finales de conférence NHL,
Ligue des Champions UEFA, et tableaux hommes et femmes de Roland
Garros à partir du R16) ; et la Coupe du Monde FIFA 2026 R16+
(15 matchs), avec prédiction émise avant le coup d'envoi du 11 juin.
Tous les tableaux sont chargés depuis la base de marchés de prédiction
PulsQuant en production, pas depuis des hypothèses-jouets équilibrées.

La construction reprend là où deux fils récents se sont arrêtés. La
note \emph{Quantum Frontiers} de 2018 \cite{quantumfrontiers_wc_2018}
esquissait un encodage en qudits pour le tableau de la Coupe du Monde
mais n'a ni publié le formalisme ni implémenté les requêtes de
vraisemblance Grover qu'elle pointait comme prolongement. Le papier
Nature/QNN football \cite{nature_qnn_football_2025} s'attaque à une
cible différente : classification d'issues match par match en machine
learning quantique, laissant le problème de l'estimation structurelle
de tableau ouvert. Côté classique, \cite{brandes2025stop} abandonnent
complètement le Monte Carlo en exploitant l'indépendance entre
sous-chemins du tableau ; nous allons à l'inverse et exploitons la
\emph{dépendance} structurelle entre matchs adjacents.

\section{Cadre théorique}
\label{sec:bg}

\subsection{Marches quantiques}

Une marche aléatoire classique sur un graphe se définit par deux
ingrédients. D'une part la position actuelle du marcheur,
c'est-à-dire sur quel nœud il se trouve. D'autre part la décision du
prochain pas, tirée typiquement par un lancer de pièce. La version
quantique de cet objet remplace ces deux registres scalaires par deux
registres quantiques. Un premier registre, dit registre de
\textbf{position}, décrit l'état de localisation sur le graphe. Un
second registre, dit registre de \textbf{pièce} (le terme consacré
dans la littérature anglo-saxonne est \emph{coin}), décrit la
direction du prochain pas. Chaque registre vit dans son propre espace
de Hilbert, qui est le formalisme standard de la mécanique quantique :
les états y sont des vecteurs complexes (les amplitudes) et les
évolutions y sont des opérateurs unitaires (les portes). La dynamique
de la marche agit sur le produit tensoriel des deux espaces, ce qui
permet à la position et à la pièce d'être intriquées entre elles. La
différence avec la marche classique tient à ce que la pièce, au lieu
de produire à chaque pas un résultat binaire indépendant, peut être
placée dans une superposition cohérente des deux directions
possibles. C'est cette superposition, propagée par intrication à
travers les pas successifs, qui produit les interférences
caractéristiques des marches quantiques
\cite{aharonov1993quantum, kempe2003quantum, portugal2018quantum}.

Pour un tableau à élimination directe vu comme un arbre binaire
équilibré de profondeur $d$ ($n = 2^d$ équipes, $2^d - 1$ matchs
internes), la marche canonique parcourt l'arbre niveau par niveau vers
la racine. La pièce, dans notre encodage, est matérialisée par l'état
du qubit associé au match courant, et son résultat de mesure détermine
lequel des deux qualifiés émerge. Les accélérations algorithmiques
exponentielles de Childs, Cleve, Deotto et Farhi
\cite{childs2003exponential} sont obtenues sur un problème de
structure différente : la traversée du \emph{welded tree}, où
l'objectif est d'atteindre un nœud cible dans un graphe dont seule la
connectivité locale est révélée à chaque étape (la marche quantique
trouve la cible en temps polynomial dans la profondeur du graphe, là
où le meilleur algorithme classique connu demande un temps
exponentiel). On reprend ici la même mécanique d'encodage par paires
\og position + pièce \fg{} intriquées, mais le problème qu'on attaque
(estimer un marginal d'issue de tableau) n'admet pas, à notre
connaissance, de speedup asymptotique du même ordre, et on n'en
revendique aucun.

\subsection{Estimation d'amplitude quantique}

Pour un état $|\psi\rangle$ qui place un poids $p$ sur un sous-espace
cible, l'échantillonnage classique demande $O(1/p\epsilon^2)$ shots
pour estimer $p$ à erreur additive $\epsilon$ près. L'algorithme
d'estimation d'amplitude de Brassard, Hoyer, Mosca et Tapp
\cite{brassard2002quantum} atteint la même précision en
$O(1/p\epsilon)$ requêtes (un speedup quadratique), que Stamatopoulos
\textit{et al.} \cite{stamatopoulos2020option} ont exploité pour le
pricing end-to-end d'options call européennes sur matériel IBM
near-term. Les raffinements ultérieurs de la QAE qui suppriment le
registre d'estimation de phase \cite{suzuki2020amplitude} réduisent le
coût en ressources de plusieurs facteurs, rendant la technique
tractable sur des dispositifs NISQ
\cite{preskill2018nisq, bharti2022nisq}. Une ligne de travail
parallèle étend la même machinerie à des speedups Monte Carlo plus
généralement \cite{montanaro2015quantum}, avec le pricing du risque de
crédit comme l'une des premières démonstrations industrielles
\cite{egger2020option}.

\subsection{Finance quantique, 2024--2026}

Le benchmark industriel récent le plus directement comparable est le
pipeline Stamatopoulos et al. \cite{stamatopoulos2020option} pour le
pricing end-to-end d'options européennes par QAE sur matériel IBM
near-term, ensuite passé à l'échelle par l'équipe IBM Quantum sur ses
machines supraconductrices à 127 qubits dans l'ère NISQ 2023--2024
\cite{bharti2022nisq}. L'estimation de ressources Goldman/IBM
\cite{goldman_ibm_2021} situe le véritable avantage quantique en
pricing de dérivés à environ 7 500 qubits logiques et une T-depth
d'environ $4{,}6 \times 10^7$, un régime encore au-delà de tout
matériel actuel mais bien en dessous des estimations historiques de la
fin des années 2010.

\section{Méthode}
\label{sec:method}

\subsection{Encodage du tableau}

Soit $T = \{t_1, \dots, t_n\}$ l'ensemble des équipes participantes
avec ratings Elo $\rho_i$, et $\mathcal{M} = \{m_0, \dots, m_{n-2}\}$
l'ensemble des matchs en ordre topologique (premier tour d'abord). À
chaque match $m \in \mathcal{M}$ on associe un qubit $q_m$. La
convention $q_m = 0$ encode \og l'équipe en position haute du match a
gagné \fg{} ; $q_m = 1$ encode \og l'équipe en position basse a
gagné \fg{}.

\paragraph{Tour 0.} Les participants sont connus dès le départ. Pour
chaque match du premier tour $m$ opposant des équipes $a$ et $b$, on
calcule $p_m = 1/(1 + 10^{(\rho_b - \rho_a)/400})$, la probabilité
logistique Elo classique, et on applique
\[
R_Y\!\left(\theta_m\right) |0\rangle, \qquad
\theta_m = 2 \arcsin\!\sqrt{1 - p_m},
\]
ce qui laisse $q_m$ dans l'état $\sqrt{p_m}\,|0\rangle +
\sqrt{1-p_m}\,|1\rangle$, de sorte qu'une mesure en base
computationnelle renvoie $|0\rangle$ avec le marginal correct.

\paragraph{Tours suivants (approximation graînée).} Pour un match $m$
au tour $r > 0$, les participants réels dépendent des résultats des
deux matchs parents $\pi^t(m), \pi^b(m)$. Un encodage fidèle
conditionnerait sur les $2^{r-1}$ configurations parentales possibles
via des rotations $r$-qubit contrôlées ; cela passe à l'échelle en
$O(2^d)$ dans la profondeur du tableau. On applique à la place une
approximation \emph{graînée} : pour chacune des 4 combinaisons de
résultats parents, on identifie l'équipe qui émergerait sous un chemin
canonique suivant les têtes de série (toujours côté haut si le bit
parent est 0, toujours côté bas si 1), on calcule la probabilité
d'affrontement sous ces équipes graines, et on applique la rotation
$R_Y$ multi-contrôlée correspondante. L'approximation graînée est
exacte pour les matchs du premier tour et le long des chemins alignés
sur les $d$ axes du tableau ; elle est biaisée vers ces appariements
canoniques pour les tours plus profonds.

\subsection{Vérification en forme close de l'approximation graînée}

À côté de $Q$ (échantillonnage Qiskit Aer) et du Monte Carlo classique
non biaisé $C$, on calcule un troisième estimateur $S$ qui énumère les
$2^{n-1}$ issues du tableau en forme close sous la \emph{même} règle
d'appariement graînée que le circuit. Le triplet $\{Q, C, S\}$ permet
alors de scinder la divergence $Q$--$C$ en deux : ce qui sépare $Q$ de
$S$ est du bruit d'échantillonnage ou une erreur d'implémentation,
tandis que ce qui sépare $S$ de $C$ est le biais de l'approximation
graînée. Tout au long des expériences, on trouve $|Q - S| < 0{,}1$\,pp,
ce qui impute l'écart à $C$ à l'approximation seule.

\section{Expériences}
\label{sec:experiments}

\subsection{Implémentation}

Toutes les expériences tournent sur Qiskit 2.4.1 + Qiskit Aer 0.13 en
mode CPU laptop. On utilise 50\,000 shots par circuit quantique et
50\,000 simulations par MC classique, avec graine fixée (42) pour la
reproductibilité. Le code source et les données brutes sont
disponibles à \url{https://github.com/pulsquant/quantum-bracket}.

\subsection{Sanity check : tableau à 8 équipes}

Un tableau à 8 équipes (7 matchs) avec ratings Elo linéairement
espacés ($\rho_A = 1850, \rho_B = 1800, \dots, \rho_H = 1500$) fournit
un cadre à faible profondeur où le biais de l'approximation graînée
doit être petit. La Figure~\ref{fig:fig1} confirme que $Q$ et $C$
s'accordent à $|\Delta|_{\max} = 1{,}48$\,pp près (limite imposée par
le bruit d'échantillonnage à 50\,k shots) sur l'ensemble des 8
champions candidats.

\begin{figure}[!htbp]
\centering
\includegraphics[width=\linewidth]{figures/fig1_smoke_test.pdf}
\caption{Sanity check : tableau synthétique à 8 équipes. Marche
quantique (bleu) et Monte Carlo classique (orange), distributions des
champions s'accordent à 1,5\,pp près sur l'ensemble des équipes.}
\label{fig:fig1}
\end{figure}

\subsection{Back-tests sur vrais tableaux : NBA, NHL, UCL, RG ATP, RG WTA}
\label{sec:multi}

On charge les vrais tableaux 2025--2026 depuis la table production
\texttt{tournament\_brackets} de PulsQuant et on construit l'objet
\texttt{Bracket} directement à partir des appariements enregistrés.
Pour les Playoffs NBA et la phase à élimination directe de la Ligue
des Champions on charge le tableau complet. Pour Roland Garros
(tableaux de 128) on coupe au R16 pour garder un compte de qubits à
15. Les ratings de force sont issus d'une table maintenue à la main
mirrorant \texttt{eloratings.net} (NBA, foot) et Tennis-Elo (ATP, WTA),
listée en annexe.

Les issues réelles fournissent les cibles ground-truth du backtest. La
Table~\ref{tab:multi_tournament} résume les marginaux de la marche
quantique et du Monte Carlo classique sur le champion, le finaliste et
les demi-finalistes réels à travers les cinq tournois.

\begin{table}[!htbp]
\centering
\caption{Back-test multi-tournois : $P(\text{outcome})$ sous la marche
quantique graînée $Q$ contre le MC classique non biaisé $C$, pour les
champions, finalistes, demi-finalistes et finalistes réels à travers
cinq tournois. Les ratios $Q/C$ en gras dépassent $1{,}5\times$ ou
tombent sous $0{,}6\times$. Tous les tournois utilisent les données de
tableau de production réelles ; les tournois de tennis sont chargés à
partir du tour des 16.}
\label{tab:multi_tournament}
\small
\begin{tabular}{llrrr}
\toprule
Tournoi & Rôle réalisé / participant & $P_Q$ & $P_C$ & $P_Q/P_C$ \\
\midrule
Playoffs NBA 2026 & SAS atteint les Finales (8\textsuperscript{e}, bat OKC en WCF) & 5,87\,\% & 3,43\,\% & $\mathbf{1{,}71\times}$ \\
                  & SAS champion (Finals menées 2--0)            & 3,46\,\% & 1,13\,\% & $\mathbf{3{,}05\times}$ \\
                  & NYK atteint les Finales (favori Est)         & 13,22\,\% & 18,75\,\% & 0,71$\times$ \\
                  & NYK champion (Finals 2--0, en cours)         & 4,83\,\% & 9,59\,\% & $\mathbf{0{,}50\times}$ \\
                  & OKC tête de série \#1 (éliminé WCF)          & 18,52\,\% & 17,73\,\% & 1,04$\times$ \\
\midrule
Stanley Cup NHL 2026 & VGK atteint la finale (sweep COL 4--0 WCF) & 12,65\,\% & 20,00\,\% & $\mathbf{0{,}63\times}$ \\
                     & CAR atteint la finale (vainqueur Est CF)   & 26,93\,\% & 23,60\,\% & 1,14$\times$ \\
                     & COL meilleur Elo Ouest (éliminé WCF)       & 18,32\,\% & 19,13\,\% & 0,96$\times$ \\
                     & OTT perdant R1 (faible Elo)                & 5,24\,\% & 2,67\,\% & $\mathbf{1{,}96\times}$ \\
\midrule
UCL 2025--26      & PSG champion (favori principal) & 10,57\,\% & 10,43\,\% & 1,01$\times$ \\
                  & ARS finaliste                    & 19,00\,\% & 13,77\,\% & 1,38$\times$ \\
                  & B04 demi-finaliste               & 10,41\,\% & 4,50\,\% & $\mathbf{2{,}32\times}$ \\
                  & RMA (éliminé en QF)              & 10,10\,\% & 14,64\,\% & 0,69$\times$ \\
\midrule
RG ATP 2026 (R16+) & Zverev finaliste (favori principal) & 4,05\,\% & 38,15\,\% & $\mathbf{0{,}11\times}$ \\
                   & Cobolli finaliste (outsider)        & 6,27\,\% & 5,74\,\% & 1,41$\times$ \\
                   & Mensik demi-finaliste (outsider)    & 5,59\,\% & 1,34\,\% & $\mathbf{4{,}17\times}$ \\
                   & Arnaldi demi-finaliste (outsider)   & 3,07\,\% & 0,69\,\% & $\mathbf{4{,}44\times}$ \\
\midrule
RG WTA 2026 (R16+) & \textbf{Andreeva championne (18 ans)} & \textbf{11,01\,\%} & \textbf{6,55\,\%} & $\mathbf{1{,}68\times}$ \\
                   & \textbf{Chwalinska finaliste (Elo 1710)} & \textbf{2,15\,\%} & \textbf{0,77\,\%} & $\mathbf{2{,}78\times}$ \\
                   & Shnaider demi-finaliste (élim. Sabalenka QF) & 3,65\,\% & 2,49\,\% & 1,46$\times$ \\
                   & Sabalenka (favorite principale, élim. QF) & 54,79\,\% & 41,08\,\% & 1,33$\times$ \\
                   & Świątek (favorite principale, élim. R16) & 3,74\,\% & 37,01\,\% & $\mathbf{0{,}10\times}$ \\
\bottomrule
\end{tabular}
\end{table}

Un schéma cohérent émerge. Les parcours profonds d'outsiders réalisés
sont amplifiés : les Spurs en finale NBA depuis la
8\textsuperscript{e} tête de série après avoir éliminé la
1\textsuperscript{re} tête de série OKC ($1{,}71\times$ sur le marginal
d'atteinte des Finales, $3{,}05\times$ sur le marginal de titre), le
titre RG femmes d'Andreeva ($1{,}68\times$), le parcours en finale de
Chwalinska avec son Tennis Elo de 1710 et son classement
114\textsuperscript{e} WTA avant le tournoi ($2{,}78\times$), et les
demi-finalistes ATP outsiders Mensik et Arnaldi ($4{,}17\times$ et
$4{,}44\times$ respectivement). Les favoris qui ont perdu tôt sont
nettement sous-pondérés en miroir, avec Świątek le cas le plus net à
$0{,}10\times$ côté WTA (éliminée au R16 par Andreeva). Du même côté,
les favoris dont l'avantage Elo est robuste face à leur tirage réel
restent à $Q/C \approx 1$ : OKC à $1{,}04\times$ sur le marginal de
champion NBA, bien qu'éliminé par SAS en WCF, simplement parce que
l'Elo d'OKC (1745) était le plus haut du tournoi et l'approximation
graînée lit son parcours dans l'Ouest comme compétitif. Quand le chalk
tient sans surprise, la marche renvoie essentiellement la valeur
classique : PSG remportant la Ligue des Champions donne
$Q/C \approx 1{,}01$, preuve que l'encodage ne fabrique pas de signal
d'upset en l'absence de celui-ci.

Le prix de cette asymétrie se voit sur les têtes de série qui
survivent profondément. Les New York Knicks ont atteint les Finales
NBA depuis l'Est, et pourtant $Q$ assigne à leur marginal d'atteinte
des Finales $P_Q = 13{,}22\,\%$ contre $P_C = 18{,}75\,\%$
($0{,}71\times$), et leur marginal de champion $P_Q = 4{,}83\,\%$
contre $P_C = 9{,}59\,\%$ ($0{,}50\times$). L'approximation graînée
pénalise les positions hautes du tableau dès lors que leur avance Elo
est modérée au point d'être érodée par des appariements conditionnels
répétés ; l'amplification des chemins outsiders et la sous-pondération
des chemins chalk sont les deux faces de la même opération. La finale
Vegas--Colorado de la conférence Ouest NHL est le contrôle le plus
propre dont on dispose : narrativement un upset, statistiquement un
pile-ou-face (Elo 1645 contre 1670), et la marche n'y produit pas non
plus de signal d'upset ($Q/C = 0{,}63\times$ sur Vegas atteignant la
finale Stanley Cup).

L'estimateur en forme close $S$ produit des valeurs à bruit
d'échantillonnage près de $Q$ sur les cinq tournois
($|Q - S|_{\max} < 0{,}09$\,pp), confirmant que la divergence $Q$--$C$
est bien le biais de l'approximation graînée et non un artefact
d'implémentation quantique.

La Figure~\ref{fig:fig2} visualise la distribution marginale du
champion des Playoffs NBA 2026 côte-à-côte pour $Q$ et $C$, illustrant
la redistribution caractéristique de la masse de probabilité depuis
les têtes de série vers les chemins du bas de tableau.

\begin{figure}[!htbp]
\centering
\includegraphics[width=\linewidth]{figures/fig2_nba_2026.pdf}
\caption{Distribution des champions des Playoffs NBA 2026 (15 qubits).
La marche quantique redistribue la masse de probabilité de manière
sélective : les têtes de série moyennes de l'Est NYK et DEN sont
sous-pondérées (leurs chemins classiques sont chargés en adversaires
de Elo modéré que l'approximation graînée lit comme vulnérables aux
upsets), tandis que les têtes de série \#1 OKC et BOS restent proches
de leurs valeurs classiques (leur avance Elo est suffisamment robuste
face au tirage réel pour qu'aucune sous-pondération ne s'applique).
SAS, 8\textsuperscript{e} tête de série de l'Ouest, a fait le parcours
profond structurellement pertinent des Playoffs 2026 en éliminant OKC
en finale de conférence Ouest pour atteindre les Finales NBA : c'est
exactement le type de chemin d'upset structurellement corrélé que la
marche quantique graînée est conçue pour amplifier ex ante. La série
des Finales NBA était en cours au moment de la rédaction (NYK menait
2--0 dans la série au meilleur des 7).}
\label{fig:fig2}
\end{figure}

\subsection{Amplification de Grover sur une requête de tableau}
\label{sec:grover}

Un prolongement naturel de l'encodage consiste à amplifier le marginal
d'une issue de tournoi spécifique au moyen de l'algorithme de Grover
\cite{grover1996fast}, plutôt que de l'estimer par échantillonnage
direct. Le tableau NBA 2026 fournit un banc d'essai propre : SAS
atteint les Finales NBA depuis la 8\textsuperscript{e} tête de série
correspond à une contrainte sur 4 qubits du registre à 15 qubits (les
bits indexant le match R1 OKC--SAS, le demi-tour Ouest CSF, la finale
de conférence Ouest et la position des Finales NBA côté Ouest). On
encode cette contrainte comme un oracle de phase $O$ et on le couple à
l'opérateur de diffusion $D = A (2|0\rangle\langle 0| - I) A^\dagger$,
où $A$ est l'unitaire de préparation du tableau de la
Section~\ref{sec:method}.

La Figure~\ref{fig:fig4} montre la courbe d'amplification empirique :
la fréquence de mesure SAS-champion, échantillonnée à $50\,000$ shots
par point, en fonction du nombre $k$ d'itérations de Grover appliquées
à $A|0\rangle$. La valeur de référence $p_0 = 3{,}46\,\%$ correspond à
la prédiction de l'encodage conditionnel à $k=0$. Les itérations
successives suivent l'enveloppe de manuel
$\sin^2((2k+1)\theta)$ avec $\sin\theta = \sqrt{p_0}$, atteignant un
pic à $97{,}89\,\%$ sur $k = 4$ (l'optimum théorique étant
$k^\star = \pi / (4 \arcsin\sqrt{p_0}) \approx 4{,}20$), redescendant à
$\sim 1\,\%$ à $k = 8$, puis remontant à $99{,}72\,\%$ sur $k = 12$.

L'implication en termes de nombre de shots est immédiate. En visant
une précision relative de $5\,\%$ sur la fréquence SAS-champion,
l'échantillonnage direct demande
$N_{\text{unamp}} = (1 - p_0)/(p_0 \rho^2) \approx 11\,200$ shots,
tandis que l'échantillonnage depuis l'état amplifié avec $k = 4$
demande $N_{\text{amp}} \approx 9$ shots (et $\approx 1$ shot à
$k = 12$). La réduction empirique d'environ $10^4 \times$ sur le
nombre de shots est cohérente avec le speedup quadratique de Grover
$O(1/\sqrt{p_0})$ au sens de l'amplification d'amplitude
\cite{brassard2002quantum}. Le coût est la profondeur du circuit :
chaque itération de Grover ajoute une application de $A$, $O$ et
$A^\dagger$, multipliant la profondeur par environ $2k + 1$. Sur Aer à
15 qubits, cela reste tractable sur un CPU monoposte, le balayage
complet $k = 0\dots15$ se terminant en quelques minutes.

L'amplification fait tourner l'état du tableau jusqu'à ce que le
sous-espace SAS-champion domine la mesure, mais ne récupère pas
directement la valeur $p_0$ elle-même : elle la fait juste tourner
près de 1. La Section~\ref{sec:qae} clôt cette lacune avec la
procédure canonique d'estimation d'amplitude Brassard-Hoyer-Mosca-Tapp
appliquée au même opérateur de Grover.

\begin{figure}[!htbp]
\centering
\includegraphics[width=0.85\linewidth]{figures/fig4_grover_speedup.pdf}
\caption{Amplification de Grover sur la requête SAS-champion du
tableau Playoffs NBA 2026. Bleu : mesure Aer à $50\,000$ shots par
point. Gris pointillé : enveloppe théorique $\sin^2((2k+1)\theta)$
avec $\sin\theta = \sqrt{p_0}$ et $p_0 = 3{,}46\,\%$ (orange en
pointillé). Deux pics à $k = 4$ ($97{,}89\,\%$) et $k = 12$
($99{,}72\,\%$), cohérents avec la périodicité de manuel de Grover.}
\label{fig:fig4}
\end{figure}

\subsection{Estimation d'amplitude quantique sur la même requête}
\label{sec:qae}

L'extension canonique de l'amplification de Grover de la section
précédente est l'estimation d'amplitude quantique (QAE)
\cite{brassard2002quantum}, qui renvoie une estimation de $p_0$
lui-même plutôt que d'amplifier le sous-espace cible jusqu'à un pic.
La construction ajoute $m$ qubits ancillaires au registre du tableau à
$n_{\text{data}} = 15$ qubits et procède en trois étapes : (i) la
préparation $A$ du tableau est appliquée au registre de données,
(ii) pour $j = 0, \dots, m - 1$, le qubit ancillaire $j$ contrôle
$Q^{2^j}$ où $Q$ est le même opérateur de Grover qu'à la
Section~\ref{sec:grover}, avec la convention de signe canonique de
sorte que ses valeurs propres sur le sous-espace bon/mauvais sont
exactement $e^{\pm 2i\theta}$ avec $\sin\theta = \sqrt{p_0}$, (iii)
une QFT inverse sur les ancillas suivie d'une mesure renvoie un entier
$y$ qui encode l'estimation d'amplitude
$\hat p = \sin^2(\pi y / 2^m)$ sur une grille à $2^m$ points.

On balaye $m \in \{4, 5, 6\}$, soit respectivement 19, 20 et 21
qubits, ce qui reste tractable sur un simulateur Aer monoposte (chaque
$m$ tourne en quelques minutes). 10\,000 shots par configuration. La
Table~\ref{tab:qae} résume la convergence vers la valeur de référence
$p_0 = 3{,}46\,\%$ établie par la baseline d'échantillonnage non
amplifié. Le mode de la mesure ancillaire se sature à
$\hat p = 3{,}81\,\%$ pour les trois valeurs de $m$ : la grille
discrète à $2^m$ points place son point le plus proche à
$y = 15 \cdot 2^{m-4}$ avec $\sin^2(\pi \cdot 15/16) \approx 3{,}81\,\%$,
ce qui se trouve dans la bande de résolution théorique $\pi / 2^m$ de
la vraie valeur pour tout $m$. L'estimateur en moyenne pondérée
$\hat p_{\text{mean}} = \sum_y P(y) \sin^2(\pi y / 2^m)$ raffine
encore avec $m$ : $3{,}84\,\%$ à $m = 4$, $3{,}90\,\%$ à $m = 5$,
$3{,}95\,\%$ à $m = 6$, tous à moins de $0{,}5$\,pp de la référence.

\begin{table}[!htbp]
\centering
\caption{Convergence de la QAE sur la requête SAS-champion. Référence
$p_0 = 3{,}46\,\%$ obtenue par échantillonnage non amplifié à
$50\,000$ shots (Figure~\ref{fig:fig4}, $k = 0$).
$\hat p_{\text{mode}}$ : mode de l'histogramme de mesure des ancillas,
décodé via $\sin^2(\pi y / 2^m)$. $\hat p_{\text{mean}}$ : estimateur
empirique pondéré par les fréquences. Le speedup quadratique de
Brassard se manifeste sous la forme des
$O(1/\sqrt{p_0}) \approx 5$ itérations $Q$ contrôlées suffisantes à
$m = 6$, contre les $O(1/p_0) \approx 30$ échantillons non amplifiés
nécessaires pour le même nombre attendu de comptes.}
\label{tab:qae}
\small
\begin{tabular}{cccccc}
\toprule
$m$ & Qubits totaux & Résolution $\pi / 2^m$ & Comptes pic $/ 10^4$ &
$\hat p_{\text{mode}}$ & $\hat p_{\text{mean}}$ \\
\midrule
4 & 19 & $0{,}196$ & $5\,037$ & $3{,}81\,\%$ & $3{,}84\,\%$ \\
5 & 20 & $0{,}098$ & $4\,943$ & $3{,}81\,\%$ & $3{,}90\,\%$ \\
6 & 21 & $0{,}049$ & $4\,792$ & $3{,}81\,\%$ & $3{,}95\,\%$ \\
\bottomrule
\end{tabular}
\end{table}

La Figure~\ref{fig:fig5} montre la distribution empirique de
$\hat p$ à $m = 6$. Le mode à $y = 60$ (replié sur $y = 4$ par la
symétrie $y \leftrightarrow 2^m - y$) porte $47{,}9\,\%$ des shots et
correspond à $\hat p = 3{,}81\,\%$ ; les pics suivants se situent à
$y = 28$ et $y = 36$, qui décodent à des amplitudes autour de
$90\,\%$ mais ne récoltent chacun que $< 1\,\%$ des shots, une queue
d'échantillonnage cohérente avec la queue de précision QPE
\cite{brassard2002quantum, suzuki2020amplitude}.

\begin{figure}[!htbp]
\centering
\includegraphics[width=0.85\linewidth]{figures/fig5_qae_estimation.pdf}
\caption{Distribution empirique de l'estimateur QAE à $m = 6$ ancillas
($21$ qubits au total, $10\,000$ shots). Bleu : histogramme des
amplitudes décodées $\hat p = \sin^2(\pi y / 2^m)$. Orange en
pointillé : référence $p_0 = 3{,}46\,\%$ depuis l'échantillonnage non
amplifié. Rouge en pointillés : estimation du mode QAE
$\hat p_{\text{mode}} = 3{,}81\,\%$.}
\label{fig:fig5}
\end{figure}

La comparaison du nombre de shots avec le schéma d'amplification de
la Section~\ref{sec:grover} clarifie la différence pratique entre les
deux méthodes. L'amplification fait tourner $p_0$ jusqu'à un pic
proche de $1$, ce qui rend bon marché la vérification de l'atteinte
du sous-espace cible, mais ne renvoie pas $p_0$ lui-même : extraire
une estimation numérique de l'état amplifié demande encore une
campagne de mesure séparée pour défaire l'angle d'amplification. La
QAE renvoie $\hat p$ directement depuis un seul shot du registre
ancillaire, au prix de l'exécution de $2^m - 1 \approx 63$ itérations
de Grover par shot à $m = 6$ contre l'application unique de $A$ par
échantillon non amplifié. L'économie asymptotique en shots est le
speedup quadratique bien connu de Brassard $O(1/p_0)$ contre
$O(1/p_0^2)$ pour une précision additive ; à $p_0 = 3{,}5\,\%$ c'est
une réduction de l'ordre de $\sim 30 \times$ sur les shots, que la
convergence empirique de la Table~\ref{tab:qae} corrobore.

\subsection{Prédiction prospective : Coupe du Monde FIFA 2026}

La Coupe du Monde FIFA 2026 démarre le 11 juin dans 16 villes hôtes
aux États-Unis, au Canada et au Mexique. On émet une prédiction
prospective sur la phase R16+ (15 matchs, 15 qubits) en prenant les
16 meilleures équipes classées par $P(\text{atteint le R16})$ depuis
le snapshot Monte Carlo classique PulsQuant du 2026-06-06T17:04, et
en les plaçant dans un tableau équilibré standard 1-contre-16. La
Table~\ref{tab:wc} montre les têtes de la distribution issue de la
marche quantique aux côtés des valeurs classiques correspondantes. La
masse se déplace des favoris de tête (FRA $-9{,}5$\,pp, ESP $-4{,}8$\,pp)
vers le milieu de hiérarchie (NED $+3{,}0$\,pp, BEL $+2{,}9$\,pp,
AUT $+2{,}2$\,pp, COL $+2{,}0$\,pp) : c'est la même redistribution que
dans le back-test NBA, transposée à un tableau dont la vérité terrain
est encore à six semaines.

\begin{table}[!htbp]
\centering
\caption{Coupe du Monde FIFA 2026, probabilité prospective de champion
sur la phase R16+ (12 entrées de tête, triées par $P_Q$). Les
appariements réels seront déterminés par la phase de groupes et
peuvent différer de notre hypothèse équilibrée 1-contre-16. Le motif
d'amplification du milieu de hiérarchie (NED, BEL, COL, AUT entre
$+2$ et $+3$\,pp) reproduit celui du backtest NBA de la
Section~\ref{sec:multi}.}
\label{tab:wc}
\begin{tabular}{lrrr}
\toprule
Équipe & $P_Q$ (marche quantique) & $P_C$ (MC classique) & $P_Q - P_C$ \\
\midrule
ARG & 31,06\,\% & 24,92\,\% & $+6{,}14$\,pp \\
BRA & 9,26\,\% & 13,14\,\% & $-3{,}88$\,pp \\
GER & 7,54\,\% & 6,04\,\% & $+1{,}49$\,pp \\
POR & 7,19\,\% & 8,36\,\% & $-1{,}17$\,pp \\
NED & 6,70\,\% & 3,68\,\% & $+3{,}02$\,pp \\
BEL & 6,11\,\% & 3,18\,\% & $+2{,}93$\,pp \\
ESP & 5,36\,\% & 10,14\,\% & $-4{,}78$\,pp \\
FRA & 5,17\,\% & 14,63\,\% & $-9{,}45$\,pp \\
COL & 4,33\,\% & 2,30\,\% & $+2{,}03$\,pp \\
ENG & 4,32\,\% & 5,97\,\% & $-1{,}65$\,pp \\
SUI & 2,80\,\% & 1,49\,\% & $+1{,}31$\,pp \\
AUT & 2,61\,\% & 0,39\,\% & $+2{,}22$\,pp \\
\bottomrule
\end{tabular}
\end{table}

\section{Discussion}
\label{sec:discussion}

\subsection{Portée de la revendication}

Un mot sur ce qu'on revendique et ce qu'on ne revendique pas. À
l'échelle simulateur 15 qubits, on ne court pas contre le MC
classique : l'échantillonnage quantique tourne en $0{,}16$~s et le
MC classique en $0{,}06$~s (Qiskit Aer sur un seul thread CPU,
50\,000 shots et simulations). La contribution est l'encodage
lui-même, pas son runtime actuel. L'approximation graînée est en
particulier une source de biais contrôlée et réversible : le check en
forme close de la Section~\ref{sec:method}.2 lui impute l'intégralité
de l'écart $Q$--$C$, ce qui signifie que monter en gamme sur
l'encodage (paragraphe suivant) récupère le comportement MC non
biaisé jusqu'au budget de qubits.

\subsection{Ce que capture l'encodage}

Ce que la construction achète, c'est la corrélation structurelle
entre matchs adjacents pour rien, par intrication, sans
conditionnement de chemin explicite. L'effet empirique sur les cinq
tournois testés est cohérent : tous les parcours profonds d'outsiders
réalisés pour lesquels on a un back-test sont amplifiés, par des
facteurs allant de $1{,}71\times$ (SAS en Finales NBA) à
$4{,}44\times$ (Arnaldi comme demi-finaliste ATP outsider), et la
championne WTA Andreeva se situe au milieu de cette plage à
$1{,}68\times$. La sous-pondération des têtes de série qui survivent
profondément, le marginal d'atteinte des Finales des New York Knicks
à $0{,}71\times$ et leur marginal de champion à $0{,}50\times$ étant
les exemples les plus nets, est le prix à payer pour cette asymétrie,
et figure dans la même Table. Les deux effets viennent de la même
approximation : pas de free lunch ici, juste un trade-off dont le
signe est exactement celui qui correspond à la manière dont les
marchés de prédiction mispricent les tournois.

\subsection{Contexte des marchés de prédiction}

L'angle marchés de prédiction est la motivation pratique. Les marchés
de prédiction en général \cite{wolfers2004prediction} et leurs
extensions combinatoires \cite{hanson2003combinatorial} sont bien
compris comme instruments de découverte de prix, mais leurs
implémentations au niveau du tableau encodent rarement les
corrélations structurelles entre matchs. Polymarket a dépassé
44~milliards de dollars de volume échangé en 2025, en route vers une
valorisation de 9~milliards adossée à ICE, et son pricing futures
agrège sur des fréquences marginales historiques plutôt que sur des
probabilités jointes structurelles. Un modèle qui assigne un poids a
priori plus élevé aux chemins d'upset structurellement corrélés
identifierait, ex ante, les mispricings côté outsider de ce profil
de liquidité, qui est aussi le côté où les spreads sont les plus
larges et l'inventaire le plus mince. C'est l'asymétrie qu'exploite
la marche graînée.

\subsection{Une question peut-être un peu bête}
\label{sec:candide}

On peut poser la question naïvement : un spectateur qui regarde le
match peut-il être quantiquement intriqué avec le ballon, l'arbitre,
ou l'attaquant au moment où celui-ci marque ? La réponse opérationnelle
est non. La décohérence macroscopique tue tout : un objet de la taille
d'un ballon perd sa cohérence quantique en environ $10^{-23}$ secondes
à température ambiante \cite{joos_zeh_1985}, et le cerveau humain est
énergétiquement un environnement encore plus hostile
\cite{tegmark_2000}. La question, prise littéralement, est close.

On remarque néanmoins, en relisant l'encodage écrit en
Section~\ref{sec:method}, qu'on retrouve mathématiquement le même type
d'objet que ceux étudiés depuis une quinzaine d'années dans le champ
encore jeune de la \emph{quantum biology} \cite{lambert2013quantum}.
Notre marche quantique sur le graphe des matchs d'un tournoi est de
la même famille que la marche excitonique du complexe FMO de la
photosynthèse \cite{engel_2007, plenio2008dephasing}, où des paquets
d'énergie cohérents se propagent sur un graphe de chromophores
conditionnellement couplés. La paire de qubits intriqués qui code une
rencontre conditionnelle dans notre construction est de la même
famille que la paire de radicaux intriqués que les oiseaux migrateurs
entretiennent dans leur cryptochrome rétinienne pour s'orienter au
champ magnétique terrestre \cite{hore_mouritsen_2016}. Les échelles
physiques sont séparées par treize ordres de grandeur, mais l'objet
mathématique sous-jacent, un état multipartite intriqué dont les
corrélations sont conditionnelles à un graphe, est rigoureusement le
même.

On ne pousse pas plus loin ici. Notre modèle reste un outil de
pricing, pas un substrat physique. Mais que la même primitive
apparaisse à la fois à l'échelle femto-seconde du transport d'énergie
dans les feuilles et à l'échelle horaire d'un tableau d'élimination
directe nous semble valoir la mention, moins comme un argument pour
ce papier que comme une indication sur ce que cette structure pourrait
pricer ailleurs.

\subsection{Extensions et travaux à venir}

L'amplification de Grover de la Section~\ref{sec:grover} et la QAE
de la Section~\ref{sec:qae} couvrent ensemble la plus simple
variante de la famille de requêtes de vraisemblance de chemin déjà
pointée par la note \emph{Quantum Frontiers} de 2018
\cite{quantumfrontiers_wc_2018} : un sous-espace inconditionnel
\og l'équipe $X$ remporte le tournoi \fg{}, interrogé à la fois par
amplification (qui fait tourner $p_0$ jusqu'à un pic) et par
estimation d'amplitude (qui renvoie $\hat p_0$ directement). L'étape
naturelle suivante est la forme conditionnelle \og sachant que $X$
atteint les demi-finales, qui est le champion le plus probable ? \fg{},
qui couple la même construction d'oracle à un registre supplémentaire
de conditionnement et qu'on laisse aux travaux à venir. Un encodage
fidèle de l'ensemble des $2^{r-1}$ configurations parentales
éliminerait en parallèle le biais graîné, au prix d'une profondeur de
circuit qui croît avec la profondeur du tableau. Les variantes QAE
NISQ-friendly sans estimation de phase \cite{suzuki2020amplitude}
réduisent substantiellement le coût en ancillas et constituent une
cible naturelle pour le matériel near-term. Toutes les extensions
restent tractables en simulation jusqu'à environ 25 qubits sur un seul
GPU via cuQuantum \cite{cuquantum_benchmark_2023}.

\FloatBarrier
\section{Conclusion}

Nous avons écrit un encodage de marche quantique pour l'estimation de
probabilités structurées par tableau, formalisé l'approximation
graînée sur laquelle il s'appuie aux tours profonds, et vérifié que
l'implémentation Qiskit Aer s'accorde, à bruit d'échantillonnage
près, avec une énumération en forme close de la même règle. À
l'échelle simulateur actuelle (15 qubits), la contribution est
conceptuelle plutôt qu'algorithmique : on ne revendique pas
d'avantage quantique, mais l'encodage reproduit une structure de
corrélation de tableau que le Monte Carlo classique ne récupère qu'au
prix d'un conditionnement de chemin explicite.

Les back-tests sur cinq vrais tableaux 2025--2026 constituent le
message empirique principal. La marche amplifie tous les parcours
profonds d'outsiders réalisés de notre échantillon (Spurs NBA,
Andreeva et Chwalinska à RG WTA, Mensik et Arnaldi à RG ATP), avec
des ratios $Q/C$ entre $1{,}71\times$ et $4{,}44\times$. Elle ne
produit pas de signal d'upset là où l'écart de rating ne le justifie
pas (PSG champion UCL à $\approx 1{,}01$ ; Vegas vs Colorado
quasi pile-ou-face à $0{,}63\times$ ; OKC champion NBA à
$1{,}04\times$). Elle sous-pondère les têtes de série dont l'avance
Elo est modérée au point d'être érodée par des appariements
conditionnels répétés (marginal d'atteinte des Finales des NYK à
$0{,}71\times$, NYK champion à $0{,}50\times$). Les trois
comportements viennent de la même approximation et sont reportés dans
la même Table. Pour le pricing sur marchés de prédiction, le signe du
biais va dans la bonne direction : le sous-pricing des chemins
d'upset structurellement corrélés est précisément la population que
l'approximation graînée identifie préférentiellement.

On implémente en outre, sur le même encodage, à la fois
l'amplification d'amplitude de Grover (Section~\ref{sec:grover}) et
l'estimation d'amplitude quantique (Section~\ref{sec:qae}) pour le
sous-espace SAS-champion du tableau NBA. La courbe d'amplification
atteint son pic à $97{,}89\,\%$ sur $k = 4$ itérations, et la QAE
renvoie $\hat p_{\text{mode}} = 3{,}81\,\%$ à $m = 6$ ancillas
(21 qubits au total), dans la bande de résolution $\pi/2^m$ de la
référence $p_0 = 3{,}46\,\%$ obtenue par échantillonnage non
amplifié. Ensemble, les deux requêtes réalisent le pipeline canonique
Brassard-Hoyer-Mosca-Tapp sur un tableau réel, refermant la plus
basique des extensions basées sur Grover initialement pointée comme
travail futur par la note \emph{Quantum Frontiers} de 2018.

Enfin nous émettons, avant le coup d'envoi du 11 juin, une prédiction
prospective pour la Coupe du Monde FIFA 2026 R16+ sous l'hypothèse
d'un tableau équilibré 1-contre-16. Les nations de milieu de
hiérarchie (NED, BEL, AUT, COL) sont amplifiées de 2 à 3 points de
pourcentage par rapport au MC classique, et les favoris de tête
(FRA, ESP, BRA) réduits d'autant. L'encodage passe à l'échelle de
l'ordre de 25 qubits sur un seul GPU via cuQuantum, et constitue un
candidat naturel pour la fenêtre NISQ tardive et pour le régime 50 à
100 qubits logiques attendu des premières générations de matériel
tolérant aux fautes.

\bibliographystyle{plain}
\bibliography{references}

\end{document}
